\section{The incidence atom compiler}
\label{sec:compiler}

Let \(I(P;\C)\) denote the locally finite incidence algebra. Its product is
\begin{equation}
(f*g)(a,b)=\sum_{a\mid c\mid b}f(a,c)g(c,b).
\label{eq:incidence-product}
\end{equation}
The zeta and Möbius elements are
\begin{equation}
\zeta(a,b)=\id_{\{a\mid b\}},
\qquad
\mu=\zeta^{-1},
\qquad
\mu(a,b)=\mob(b/a).
\label{eq:zeta-mobius}
\end{equation}
For every source object \(n\), including composites, introduce the uniform
coordinate idempotent
\begin{equation}
\epscoord_n(a,b)=\id_{\{a=b=n\}}.
\label{eq:coordinate-idempotent}
\end{equation}

\begin{definition}[Incidence compiler]
\label{def:compiler}
The compiled source coordinate is
\begin{equation}
q_n=\zeta*\epscoord_n*\mu.
\label{eq:compiler}
\end{equation}
\end{definition}

No factorization or primality test is used in \cref{def:compiler}. Directly
expanding the two convolutions gives
\begin{equation}
q_n(a,b)
=\id_{\{a\mid n\mid b\}}\mob(b/n).
\label{eq:kernel}
\end{equation}
This formula will later yield the countable rank-one representation.

\begin{theorem}[Complete primitive-idempotent system]
\label{thm:primitive-system}
For every finite cutoff \(P_N\), and intervalwise in the countable incidence
algebra,
\begin{equation}
q_nq_m=\delta_{nm}q_n,
\qquad
\sum_n q_n=\delta.
\label{eq:complete-idempotents}
\end{equation}
Each \(q_n\) is primitive.
\end{theorem}

\begin{proof}
Using \(\mu\zeta=\delta\),
\[
q_nq_m
=\zeta\epscoord_n\mu\zeta\epscoord_m\mu
=\zeta\epscoord_n\epscoord_m\mu
=\delta_{nm}q_n.
\]
At finite cutoff, \(\sum_n \epscoord_n=\delta\), whence
\[
\sum_n q_n
=\zeta\left(\sum_n \epscoord_n\right)\mu
=\zeta\delta\mu=\delta.
\]
The countable sum has only finitely many contributions on every interval.
Each \(\epscoord_n\) is primitive in the finite incidence algebra and
conjugation by the unit \(\zeta\) preserves primitivity. The countable claim
is the compatible intervalwise statement.
\end{proof}

The compiler itself does not select primes: it produces one idempotent for
every integer. Selection enters through the fixed source predicate.

\begin{definition}[Source-derived letter action]
\label{def:letter-action}
For a positive-integer return letter \(n\), define
\begin{equation}
A_n=\id_{\At(P)}(n)q_n,
\label{eq:letter-action}
\end{equation}
where membership in \(\At(P)\) is evaluated by
\eqref{eq:atom-predicate}.
\end{definition}

\begin{theorem}[Exact necklace-resolved atom selector]
\label{thm:word-selector}
For every nonempty word \(w=n_1\cdots n_r\),
\begin{equation}
\Tr(A_{n_1}\cdots A_{n_r})
=
\begin{cases}
1,&n_1=\cdots=n_r=p\in\At(P),\\
0,&\text{otherwise}.
\end{cases}
\label{eq:word-selector}
\end{equation}
The coefficient is invariant under cyclic rotation. The selected primitive
cyclic classes are precisely the one-letter source-atom loops, and their
\(r\)-fold temporal repetitions have weight
\begin{equation}
u^{r\ell(p)}p^{-rs}
\label{eq:repetition-weight}
\end{equation}
with multiplicity one.
\end{theorem}

\begin{proof}
If a letter is not an atom, its action is zero. If all letters are atoms but
two differ, \cref{thm:primitive-system} makes the product zero. A
monochromatic atom word equals \(q_p^r=q_p\). At a finite cutoff \(q_p\) is
similar to the rank-one coordinate idempotent and has trace one. The
countable trace-one statement is proved directly in \cref{thm:rank-one}.
Cyclic invariance follows either from the trace or directly from the
dichotomy.
\end{proof}

\begin{remark}[Source exponent versus time]
\label{rem:source-time}
The single source letter \(p^r\) is a composite integer and is killed by
\eqref{eq:letter-action}. The word with \(r\) copies of \(p\) is the temporal
repetition of the primitive loop and survives. The exponent in
\eqref{eq:repetition-weight} records both distinctions: \(p^{-rs}\) comes
from \(r\) traversals, and \(u^{r\ell(p)}\) records all underlying digit
steps.
\end{remark}

This is the analytic A1 advance. It is stronger than an equality obtained
only after commuting letter variables or aggregating endpoints: mixed and
composite-letter products vanish before the cyclic trace is taken.
